% Version 4.2.1 end matter: one full-page figure per page after the report text.
\clearpage
\onecolumn

\begin{figure}[p]
\centering
\includegraphics[width=\textwidth,height=.84\textheight,keepaspectratio]{fig1_bands_quantum_geometry.pdf}
\caption{\textbf{Continuum bands and quantum geometry.} (a,b) The first ten moir\'e bands along $\Gamma$--$K$--$M$--$\Gamma$ at $\theta=2.6^\circ$ and $3.0^\circ$; bands 1--3 are colored blue, orange and green, respectively. (c--e) Berry curvature and (f--h) trace quantum metric for bands 1--3 at $\theta=2.6^\circ$ across the hexagonal Brillouin zone. Both rows use a common physical color scale with explicit numeric ticks.}
\label{fig:bands}
\end{figure}

\clearpage
\begin{figure}[p]
\centering
\includegraphics[width=\textwidth,height=.84\textheight,keepaspectratio]{fig2_multiband_ed.pdf}
\caption{\textbf{Multiband ED on the $3\times3$ cluster.} Top row: $\nu=1/3$; bottom row: $\nu=2/3$. (a,d) Ground-state energy per hole versus retained bands. (b,e) The lowest three levels in every momentum sector for one-, three- and five-band ED, plotted as $E-E_{\rm GS}$ in meV. Larger opaque markers denote sector ground states; the next two levels are smaller and semitransparent. (c,f) Normalized band-resolved ground-state weights. Charge-$U(1)$, spin-$U(1)$ and translations are exact; $C_3/C_6$ equivalence accelerates the largest five-band sectors.}
\label{fig:multiband}
\end{figure}

\clearpage
\begin{figure}[p]
\centering
\includegraphics[width=\textwidth,height=.84\textheight,keepaspectratio]{fig3_fig5_spectra_observables.pdf}
\caption{\textbf{27-site one-band reproduction of Supplementary Fig.~5.} Top row: $\nu=1/3$; bottom row: $\nu=2/3$. (a,d) Momentum-resolved spectra, plotted as $E-E_{\rm GS}$ in meV. Red markers denote the CDW or FCI three-state manifold. (b,e) Ground-manifold averaged $n(\bm k)$ over weak, interpolated band-1 Berry curvature. (c,f) Projected $S(\bm q)$ over the identical weak, interpolated band-1 $\mathrm{Tr}\,g_1-|\Omega_1|$ background. Bright marker colorbars encode the observables independently.}
\label{fig:fig5}
\end{figure}

\clearpage
\begin{figure}[p]
\centering
\includegraphics[width=\textwidth,height=.84\textheight,keepaspectratio]{fig4_manybody_berry_curvature.pdf}
\caption{\textbf{Twisted-boundary ED and many-body topology on the 27-cell one-band cluster.} Top row: $\nu=1/3$; bottom row: $\nu=2/3$. (a,d) Spectral flow versus $\phi_1$ at $\phi_2=0$. (b,e) Spectral flow versus $\phi_2$ at $\phi_1=0$. The lowest three energies are red. Gray points show additional sampled levels. (c,f) Non-Abelian curvature $\mathcal F_{\phi_1\phi_2}$ on the original $6\times6$ primitive twist grid; white dots mark FHS plaquette centers. The raw plaquette-flux sums give Chern number zero for the CDW and two for the tracked FCI manifold.}
\label{fig:manybodychern}
\end{figure}

% Figure 5 was retired in the v4 restructuring. Preserve the established
% numbering of the nine-cell and 27-cell QNS summary figures.
\setcounter{figure}{5}

\begin{figure}[p]
\centering
\includegraphics[width=\textwidth,height=.84\textheight,keepaspectratio]{fig6_9cell_qns_summary.pdf}
\caption{\textbf{Nine-cell QNS summary.} The former optimization and timing
figures are combined on one full page. Rows one and two show $\nu=1/3$ and
$2/3$ optimization traces, one--five-band ED energy comparisons and bare-band
weights. Rows three and four show all outer-$C_3$ sector traces, final-ten
sector offsets, and the four-method timing comparison. The outer timing is
the complete three-sector wall time divided by three. Error bars and shaded
energy widths are final-ten RMS fluctuations, not independent uncertainties.
The selected outer-state band bars are yellow. CG stability is omitted here
and collected in Fig.~\ref{fig:cgstability}.}
\label{fig:qns9summary}
\end{figure}

\clearpage
\begin{figure}[p]
\centering
\includegraphics[width=\textwidth,height=.84\textheight,keepaspectratio]{fig7_27cell_qns_summary.pdf}
\caption{\textbf{Twenty-seven-cell QNS summary.} The layout follows
Fig.~\ref{fig:qns9summary}.  Since five-band 27-cell ED is unavailable, the
energy panels use full $M$, $P_\Gamma[M]$ and $P_mP_\Gamma[M]$ as categorical
abscissa labels; the one-band 27-cell ED energy is a dashed horizontal
reference. All three outer-energy markers are yellow, with circle, square and
triangle denoting $m=0,1,2$; the selected outer-state band bars use the same
yellow. Rows three and four resolve the three outer sectors and report measured
training times. The hatched outer cost is the complete two-GPU
critical path divided by the three resulting candidate states.}
\label{fig:qns27summary}
\end{figure}

\clearpage
\begin{figure}[p]
\centering
\includegraphics[width=\textwidth,height=.84\textheight,keepaspectratio]{fig8_27cell_qns_observables.pdf}
\caption{\textbf{Twenty-seven-cell full-$M$ and $P_\Gamma[M]$ observables.}
Rows compare both ansatzes at $\nu=1/3$ and $2/3$. Columns show the
Bloch-projected occupation summed over bands 1--5, the band-1 occupation
conditionally normalized to $N_e$, the full static structure factor
transformed from the real-space density-pair correlation, and folded charge
density. Both momentum panels and the main $S(\bm q)$ panel use exactly the 27
physical ED momentum vectors. A directly evaluated 37-vector, $C_6$-closed
grid is retained only for boundary rotation audits. No observable is $C_3$
averaged, and no ED momentum profile enters the first-band normalization. All
values use the saved final 4,128-walker training ensembles without a separate
validation chain.}
\label{fig:nbobservables-baseline}
\end{figure}

\clearpage
\begin{figure}[p]
\centering
\includegraphics[width=\textwidth,height=.84\textheight,keepaspectratio]{fig9_27cell_qns_observables_nu1of3.pdf}
\caption{\textbf{Twenty-seven-cell $\nu=1/3$ outer-sector observables.}
Rows are the complete-wavefunction projectors $P_0,P_1,P_2$; the lowest
final-ten branch is marked directly. Columns and all momentum-observable
rendering conventions match Fig.~\ref{fig:nbobservables-baseline} and the ED
reference in Fig.~\ref{fig:fig5}. Shared color limits permit a direct visual
comparison without sector-by-sector rescaling.}
\label{fig:nbobservables13}
\end{figure}

\clearpage
\begin{figure}[p]
\centering
\includegraphics[width=\textwidth,height=.84\textheight,keepaspectratio]{fig10_27cell_qns_observables_nu2of3.pdf}
\caption{\textbf{Twenty-seven-cell $\nu=2/3$ outer-sector observables.}
Rows are $P_0,P_1,P_2$ and columns use the same definitions, color limits and
marker rules as Fig.~\ref{fig:nbobservables13}. The nearly
degenerate final-ten $m=1,2$ energies are not by themselves evidence of a
converged three-state manifold.}
\label{fig:nbobservables23}
\end{figure}

\clearpage
\begin{figure}[p]
\centering
\includegraphics[width=\textwidth,height=.72\textheight,keepaspectratio]{fig11_cg_stability_9cell_27cell.pdf}
\caption{\textbf{Supplementary nine/27-cell natural-gradient stability.}
True relative residuals are recomputed from $F-(S+\lambda I)\Delta\theta$.
Columns compare nine and 27 cells and rows compare $\nu=1/3$ and $2/3$.
Horizontal lines mark the 0.02 early-stop and 0.05 cap-acceptance thresholds;
the vertical line marks update 20. The 27-cell residual excursions and rejected
updates show that the 120-step energies remain finite-optimization diagnostics.}
\label{fig:cgstability}
\end{figure}
